\section{Introduction}

The recent surge in the capabilities of large language models (LLMs) has dramatically transformed the landscape of text generation. State-of-the-art AI systems can now produce highly coherent and contextually relevant narratives, reports, and articles that are increasingly difficult to distinguish from human-written content. While these advances unlock significant opportunities for innovation across domains such as journalism, education, and digital assistants, they also intensify concerns regarding authenticity, misinformation, and the potential misuse of synthetic text.

The problem of online fake new \cite{}, especially related to elections and diseases \cite{} is exacerbated by AI generated fake news. 

Amidst this evolving environment, the challenge of reliably detecting and attributing AI-generated text has become a research priority \cite{zellers2020defendingneuralfakenews,mitchell2023detectgptzeroshotmachinegeneratedtext}. The capability to differentiate between human-authored and machine-generated content is crucial for safeguarding the integrity of news media, academic publishing, and online discourse. However, progress in this area has been hindered by the scarcity of large, diverse, and well-annotated datasets that contain both human and AI-generated texts under realistic, real-world conditions.

In this paper, we introduce a dataset to detect AI generated text. The dataset is a comprehensive, extensively annotated resource curated to advance research on the detection of AI-generated content. Our dataset is built upon a large-scale collection of authentic news articles from the New York Times, spanning over two decades and covering a broad spectrum of topics. For each sampled article, we provide both the original human-authored narrative and multiple synthetic counterparts generated by leading LLMs, including Gemma-2-9b \cite{gemmateam2024gemma2improvingopen}, Mistral-7B\cite{jiang2023mistral7b}, Qwen-2-72B\cite{bai2023qwentechnicalreport}, LLaMA-8B \cite{touvron2023llamaopenefficientfoundation}, Yi-Large\cite{ai2025yiopenfoundationmodels}, and GPT-4-o\cite{openai2024gpt4technicalreport}. Each sample is enriched with detailed metadata, facilitating nuanced analysis and enabling a wide range of research applications.

The unique structure of the dataset supports tasks such as binary classification (distinguishing human from AI content) and model attribution (identifying the specific LLM responsible for a synthetic text). By bridging the gap between real-world journalistic data and diverse state-of-the-art generative models, our work lays the foundation for more robust benchmarks, and tools that are urgently needed to maintain trust, clarity, and accountability in the age of generative AI.

